\section{Tối ưu bằng QAT}

\subsection{Thiết lập quantization}

QAT được cấu hình cho weight và activation theo scheme \cbs{per-tensor hoặc
per-channel, symmetric hoặc asymmetric}. Những toán tử chưa hỗ trợ INT8 hoặc
nhạy với lượng tử được giữ ở precision cao hơn. Cấu hình quantization phải tương
thích với đường export ONNX/TensorRT đã lựa chọn.

\subsection{Fake Quantization}

Các fake-quant node được chèn vào graph để mô phỏng clipping và làm tròn. Giai
đoạn đầu thu thập thống kê activation; sau đó observer được đóng băng để scale ổn
định trong phần cuối huấn luyện. Báo cáo lưu histogram hoặc min--max của các
tầng nhạy cảm khi cần giải thích suy giảm accuracy.

\subsection{QAT Training}

QAT khởi tạo từ baseline hoặc checkpoint đã pruning và fine-tuning. Learning rate
được đặt nhỏ hơn huấn luyện ban đầu; BatchNorm và observer được đóng băng theo
lịch \cbs{lịch thực tế}. Quá trình theo dõi đồng thời loss, validation mAP và dấu
hiệu mất ổn định số.

\subsection{Export INT8}

Checkpoint QAT được export sang ONNX với thông tin lượng tử tương ứng. Graph
được kiểm tra bằng ONNX checker và chạy so sánh đầu ra trên một tập mẫu. Sau đó,
TensorRT build engine INT8 hoặc mixed precision; log build và danh sách layer
precision được lưu cùng artifact.

\subsection{Đánh giá mô hình QAT}

Mô hình QAT được đánh giá ở ba điểm: checkpoint fake-quant trong framework,
ONNX sau export và TensorRT engine trên thiết bị. Chênh lệch giữa các điểm giúp
xác định sai số đến từ lượng tử, export hay backend. Các chỉ số cần điền gồm
Precision, Recall, $mAP@50$, $mAP@50{:}95$, latency, FPS và kích thước engine.